\documentclass[12pt,a4paper]{article}
\usepackage[spanish]{babel}
\usepackage[utf8]{inputenc}
\usepackage[T1]{fontenc}
\usepackage{geometry}
\usepackage{hyperref}
\usepackage{longtable}
\usepackage{booktabs}
\usepackage{enumitem}
\geometry{margin=2.5cm}

\title{Informe del Estado Actual del Sistema\\Proyecto GrupoCordillera}
\author{Equipo de Desarrollo}
\date{\today}

\begin{document}
\maketitle

\section{Resumen Ejecutivo}
Este informe describe el estado actual del desarrollo del repositorio \texttt{GrupoCordillera}, con foco en la arquitectura de microservicios, la aplicacion del \textit{Repository Pattern}, la estrategia de ramas Git, la orquestacion con Docker Compose y la automatizacion CI con Jenkins.

Repositorio: \url{https://github.com/M4gic1540/GrupoCordillera.git}

\section{Arquitectura General del Proyecto}
El sistema esta organizado como un conjunto de microservicios desplegables de forma independiente, con una capa de gateway para ruteo y controles de seguridad perimetral.

\subsection{Componentes principales}
\begin{itemize}[leftmargin=*]
    \item \textbf{authservice}: Servicio de autenticacion/autorizacion (registro, login, refresh, validacion JWT y perfil de usuario).
    \item \textbf{data-ingestion-service}: Servicio de ingesta de eventos y coordinacion con el motor KPI.
    \item \textbf{kpi-engine}: Servicio de calculo y exposicion de indicadores (KPIs).
    \item \textbf{gateway}: Reverse proxy NGINX con balanceo, control de acceso, validacion de token y reglas de defensa basica.
\end{itemize}

\subsection{Despliegue e integracion}
La composicion del sistema en desarrollo se realiza con \texttt{docker-compose.yml}, incluyendo:
\begin{itemize}[leftmargin=*]
    \item 3 bases de datos PostgreSQL separadas por dominio (auth, events, kpi).
    \item 2 replicas de \texttt{authservice} y 2 replicas de \texttt{data-ingestion-service}.
    \item 1 instancia de \texttt{kpi-engine}.
    \item 1 gateway NGINX como punto de entrada.
\end{itemize}

Esta distribucion sigue el principio de aislamiento por contexto de negocio y facilita escalamiento horizontal en servicios con mayor carga.

\section{Repository Pattern y su Aplicacion}
\subsection{Concepto}
El \textit{Repository Pattern} abstrae el acceso a datos para desacoplar la logica de negocio de los detalles de persistencia. En este proyecto se implementa con Spring Data JPA, usando interfaces \texttt{JpaRepository}.

\subsection{Aplicacion en authservice}
En \texttt{authservice} se observa una aplicacion madura del patron:
\begin{itemize}[leftmargin=*]
    \item \texttt{UserRepository}: operaciones de consulta de usuario por email y verificacion de existencia.
    \item \texttt{RefreshTokenRepository}: busqueda de tokens vigentes y operaciones de limpieza/rotacion de refresh tokens.
\end{itemize}

Beneficio: el servicio de autenticacion concentra reglas de negocio (registro, login, refresh, perfil) sin mezclar SQL ni detalles de infraestructura.

\subsection{Aplicacion en data-ingestion-service}
\texttt{IngestionEventRepository} extiende \texttt{JpaRepository} y permite persistir eventos de ingesta con un contrato simple.

Beneficio: el servicio de ingesta mantiene un flujo limpio para evolucionar hacia consultas especializadas sin romper la capa de servicio.

\subsection{Aplicacion en kpi-engine}
\begin{itemize}[leftmargin=*]
    \item \texttt{KpiDefinitionRepository}: busqueda por codigo funcional de KPI.
    \item \texttt{KpiSnapshotRepository}: recuperacion de snapshots recientes para consulta.
\end{itemize}

Beneficio: separa definiciones de KPI y resultados calculados, mejorando trazabilidad y mantenibilidad.

\subsection{Valor de la decision}
La eleccion de Repository Pattern en los tres servicios permite:
\begin{itemize}[leftmargin=*]
    \item Consistencia de arquitectura entre equipos y modulos.
    \item Facilidad para pruebas unitarias/integracion.
    \item Menor acoplamiento con la base de datos y mayor portabilidad.
\end{itemize}

\section{Distribucion del Repositorio y Estrategia de Ramas}
Actualmente se utilizan 3 ramas principales:
\begin{itemize}[leftmargin=*]
    \item \textbf{main}: rama estable de referencia para entregas consolidadas.
    \item \textbf{developer}: rama de integracion continua del equipo, donde se consolidan cambios validados.
    \item \textbf{feature}: rama de trabajo para implementaciones incrementales, correcciones y experimentacion controlada.
\end{itemize}

\subsection{Por que esta decision es correcta}
\begin{itemize}[leftmargin=*]
    \item Reduce riesgo en \texttt{main}, al evitar integrar cambios no verificados.
    \item Permite validar en \texttt{developer} la compatibilidad entre microservicios antes de liberar.
    \item Facilita desarrollo paralelo en \texttt{feature} sin bloquear al equipo.
\end{itemize}

Esta estrategia es coherente con practicas de CI/CD y con la necesidad de controlar regresiones en una arquitectura distribuida.

\section{Gateway (NGINX) y Control de Trafico}
El gateway implementa funciones clave:
\begin{itemize}[leftmargin=*]
    \item Balanceo con \texttt{least\_conn} para auth e ingestion.
    \item \texttt{auth\_request} contra \texttt{/api/auth/validate} para proteger rutas privadas.
    \item Limitacion de tasa (rate limiting) por IP.
    \item Reglas WAF basicas para bloquear patrones de inyeccion y exploracion maliciosa.
    \item Propagacion de \texttt{X-Correlation-ID} para observabilidad transversal.
\end{itemize}

Esta capa mejora seguridad, resiliencia y trazabilidad operacional del sistema.

\section{Docker Compose: Estado y Justificacion}
\subsection{Aspectos positivos del disenio actual}
\begin{itemize}[leftmargin=*]
    \item Separacion de bases de datos por dominio (principio de aislamiento de datos).
    \item Dependencias con \texttt{healthcheck} para controlar orden de arranque.
    \item Configuracion por variables de entorno para portabilidad.
    \item Volumenes persistentes para no perder datos en reinicios.
\end{itemize}

\subsection{Implicaciones tecnicas}
\begin{itemize}[leftmargin=*]
    \item Se facilita el testing de integracion en entorno local reproducible.
    \item El escalamiento de servicios con replicas permite simular carga real.
\end{itemize}

\section{Jenkins y Pipeline CI}
El pipeline actual en Jenkins cubre:
\begin{itemize}[leftmargin=*]
    \item Checkout del codigo.
    \item Build \& Test por microservicio (authservice, data-ingestion-service, kpi-engine).
    \item Construccion de imagenes Docker por servicio.
    \item Publicacion de reportes de pruebas (Surefire XML).
\end{itemize}

Adicionalmente:
\begin{itemize}[leftmargin=*]
    \item Etapa SonarQube presente pero desactivada temporalmente por compatibilidad/conectividad.
    \item Etapa de despliegue stack desactivada temporalmente por ajustes de montaje/rutas.
\end{itemize}

Esto mantiene el pipeline util y estable mientras se corrigen temas de entorno, sin detener la validacion automatica de calidad basica.

\section{Estado de Pruebas}
En el estado actual se incrementaron pruebas en \texttt{authservice}, superando 20 tests y validando:
\begin{itemize}[leftmargin=*]
    \item Reglas de negocio de autenticacion (registro, login, refresh, perfil).
    \item Casos de error y validacion de entrada en controladores.
    \item Flujo de respuesta HTTP y contratos basicos de API.
\end{itemize}

Esto fortalece la confianza en cambios futuros y reduce riesgo de regresiones en CI.

\section{Paso 4: Evaluacion Etica y Recomendaciones}
\subsection{Privacidad de datos de usuarios}
\textbf{Como se garantiza actualmente:}
\begin{itemize}[leftmargin=*]
    \item Uso de autenticacion JWT para acceso controlado.
    \item Endpoints de negocio protegidos por validacion de token desde gateway.
    \item Contraseñas almacenadas con hash en \texttt{authservice} (no en texto plano).
    \item Separacion de datos por servicio y base de datos (menor superficie de exposicion cruzada).
\end{itemize}

\textbf{Riesgos a seguir gestionando:}
\begin{itemize}[leftmargin=*]
    \item Posible filtrado de informacion sensible en logs si no se enmascara.
    \item Manejo de secretos por variables de entorno sin un gestor dedicado.
\end{itemize}

\subsection{Sostenibilidad computacional}
\textbf{Aspectos favorables:}
\begin{itemize}[leftmargin=*]
    \item Escalamiento selectivo por servicio (se replica solo donde se necesita).
    \item Balanceo que evita saturacion desigual entre instancias.
    \item Separacion de responsabilidades que permite optimizacion por componente.
\end{itemize}

\textbf{Aspectos a optimizar:}
\begin{itemize}[leftmargin=*]
    \item Falta de limites explicitos de CPU/RAM en \texttt{docker-compose}.
    \item Pipeline aun sin analisis estatico continuo (SonarQube deshabilitado).
\end{itemize}

\subsection{Recomendaciones concretas de mejora}
\begin{enumerate}[leftmargin=*]
    \item \textbf{Gestion de secretos}: migrar credenciales a un vault (por ejemplo, HashiCorp Vault o secretos del orquestador).
    \item \textbf{Hardening de logs}: enmascarar tokens/correos sensibles y definir politicas de retencion.
    \item \textbf{SonarQube estable}: fijar version de plugin y Java en Jenkins para reactivar analisis automatico.
    \item \textbf{Recursos en contenedores}: definir \texttt{limits/reservations} de CPU y memoria para eficiencia energetica.
    \item \textbf{Mayor cobertura}: extender tests en \texttt{kpi-engine} y \texttt{data-ingestion-service} al mismo nivel de \texttt{authservice}.
    \item \textbf{Observabilidad}: integrar metricas y trazas (Prometheus + Grafana + trazabilidad distribuida) para detectar degradaciones tempranas.
    \item \textbf{Politica de ramas}: formalizar reglas de merge (PR obligatorio, checks de tests y aprobacion de revisores).
\end{enumerate}

\section{Matriz de Cumplimiento de Requerimientos}
\subsection{Resumen comparativo y porcentaje de cumplimiento}
\begin{longtable}{lcccc}
\toprule
\textbf{Tipo} & \textbf{Total} & \textbf{Completos} & \textbf{Parciales} & \textbf{\% Cumplimiento} \\
\midrule
RF & 18 & 6 & 7 & 33.3\% \\
RNF & 18 & 12 & 4 & 66.7\% \\
\bottomrule
\end{longtable}

\noindent
\textbf{Leyenda:}
\begin{itemize}[leftmargin=*]
    \item \textbf{Total}: cantidad de requerimientos definidos.
    \item \textbf{Completos}: implementados y verificables en código/configuración/pruebas.
    \item \textbf{Parciales}: implementados parcialmente o sin evidencia operativa suficiente.
    \item \textbf{\% Cumplimiento}: proporción de completos sobre el total.
\end{itemize}

\subsection{Leyenda de estado}
\begin{itemize}[leftmargin=*]
    \item \textbf{Completo}: Implementado y verificable en codigo/configuracion.
    \item \textbf{Parcial}: Implementado de forma incompleta o sin evidencia operativa suficiente.
    \item \textbf{Pendiente}: No implementado en el estado actual.
\end{itemize}

\subsection{Requerimientos Funcionales (RF)}
\renewcommand{\arraystretch}{1.2}
\begin{longtable}{p{1.5cm}p{6.4cm}p{2.3cm}p{3.6cm}}
\toprule
\textbf{ID} & \textbf{Descripcion} & \textbf{Estado} & \textbf{Evidencia / Brecha} \\
\midrule
\endfirsthead
\toprule
\textbf{ID} & \textbf{Descripcion} & \textbf{Estado} & \textbf{Evidencia / Brecha} \\
\midrule
\endhead
RF-01 & Autenticacion por email y contrasena. & Completo & Endpoints de registro/login implementados en authservice. \\
RF-02 & Generacion y validacion de JWT. & Completo & Servicio JWT y endpoint de validacion usados por gateway. \\
RF-03 & Gestion de roles y permisos (RBAC). & Parcial & Existen roles USER/ADMIN; falta autorizacion fina por rol en rutas de negocio. \\
RF-04 & Integracion con POS, ERP, e-commerce, CRM. & Parcial & Hay conectores externos configurables; no se evidencia cobertura explicita de POS y e-commerce. \\
RF-05 & Extraccion automatica de datos fuente. & Completo & IngestionService consume lotes via conectores por sistema fuente. \\
RF-06 & Transformacion a esquema unificado. & Parcial & Se normaliza a IngestionEvent; falta mapeo semantico avanzado por tipo de fuente. \\
RF-07 & Almacenamiento consolidado por servicio. & Completo & Bases de datos separadas por microservicio en Docker Compose. \\
RF-08 & Definicion de KPI. & Completo & KpiEngine crea/asegura definiciones (KpiDefinition). \\
RF-09 & Calculo automatico de KPI. & Completo & Recalculo automatico de snapshots en KpiEngineService. \\
RF-10 & Deteccion de desviaciones por umbral. & Pendiente & No hay logica de umbrales ni evaluacion de desviacion. \\
RF-11 & Generacion de alertas por umbral excedido. & Pendiente & No existe modulo de alertas ni notificacion activa. \\
RF-12 & API KPI filtrada por periodo/sucursal/area. & Pendiente & Solo existe consulta de ultimos snapshots sin filtros solicitados. \\
RF-13 & Reportes en PDF y Excel. & Pendiente & No hay generacion de documentos PDF/XLSX en servicios actuales. \\
RF-14 & Descarga de reportes por URL temporal. & Pendiente & No existe mecanismo de URL firmada o expirable para descarga. \\
RF-15 & Historial de reportes generados. & Pendiente & No hay entidad/repositorio de historial de reportes. \\
RF-16 & Consulta en near real-time. & Completo & Flujo de ingesta y recalc de KPI es inmediato al procesamiento de eventos. \\
RF-17 & Notificacion al motor KPI con nuevos datos. & Completo & KpiNotificationService invoca /api/kpi/recalculate tras ingesta. \\
RF-18 & Validacion de integridad antes de persistir. & Parcial & Hay validaciones de formato/flujo; faltan reglas completas de integridad de negocio y calidad. \\
\bottomrule
\end{longtable}

\subsection{Requerimientos No Funcionales (RNF)}
\begin{longtable}{p{1.5cm}p{6.4cm}p{2.3cm}p{3.6cm}}
\toprule
\textbf{ID} & \textbf{Descripcion} & \textbf{Estado} & \textbf{Evidencia / Brecha} \\
\midrule
\endfirsthead
\toprule
\textbf{ID} & \textbf{Descripcion} & \textbf{Estado} & \textbf{Evidencia / Brecha} \\
\midrule
\endhead
RNF-01 & Escalabilidad horizontal con Docker. & Completo & Replicas de authservice e ingestion + gateway balanceando trafico. \\
RNF-02 & Alta disponibilidad de servicios criticos. & Parcial & HA parcial por replicas; kpi-engine aun en instancia unica. \\
RNF-03 & Tiempo de respuesta menor a 2s. & Pendiente & No hay benchmark/SLA medido ni evidencia de pruebas de rendimiento formales. \\
RNF-04 & Soporte de concurrencia sin degradacion significativa. & Parcial & Arquitectura preparada, pero faltan pruebas de carga concurrente documentadas. \\
RNF-05 & Autenticacion/autorizacion segura con JWT. & Completo & JWT + validacion en gateway + seguridad stateless. \\
RNF-06 & Hash de contrasenas con BCrypt. & Completo & PasswordEncoder BCrypt configurado en SecurityConfig. \\
RNF-07 & Aislamiento de bases por microservicio. & Completo & postgres auth/events/kpi separados en Compose. \\
RNF-08 & Tolerancia a fallos con Circuit Breaker. & Completo & Resilience4j configurado en conectores de ingesta. \\
RNF-09 & Logs estructurados para monitoreo. & Completo & Logback + patron consistente y niveles definidos. \\
RNF-10 & Trazabilidad por Correlation ID. & Completo & Filtro de correlacion + propagacion en NGINX. \\
RNF-11 & Metricas de salud/rendimiento con Actuator. & Completo & Dependencias Actuator habilitadas en los servicios principales. \\
RNF-12 & Despliegue local con Docker Compose. & Completo & Stack ejecutable con docker-compose.yml. \\
RNF-13 & Mantenibilidad y separacion de responsabilidades. & Completo & Arquitectura por capas y microservicios desacoplados. \\
RNF-14 & Extensibilidad de conectores (Factory Method). & Completo & ConnectorFactory + SourceConnector para agregar nuevas fuentes. \\
RNF-15 & Consistencia eventual en procesamiento. & Parcial & Flujo desacoplado de ingesta/recalculo, pero notificacion actual es sincrona HTTP. \\
RNF-16 & Manejo controlado de errores. & Completo & Exception handlers y tratamiento de fallos en conectores. \\
RNF-17 & HTTPS entre servicios en produccion. & Pendiente & No hay configuracion TLS/mTLS inter-servicio en estado actual. \\
RNF-18 & Integracion futura con Prometheus/Grafana. & Completo & Base lista con Actuator y diseño compatible para observabilidad. \\
\bottomrule
\end{longtable}

\section{Conclusiones}
El proyecto presenta una base arquitectonica solida para un sistema distribuido: microservicios desacoplados, patron repositorio consistente, gateway con controles de seguridad, CI funcional y contenedorizacion reproducible.

Las decisiones tomadas (3 ramas, pruebas automatizadas, separacion por dominios, compose con healthchecks) son tecnicamente justificables y alineadas con buenas practicas de ingenieria de software.

La siguiente madurez del sistema debe enfocarse en seguridad operativa (secretos/logs), calidad continua (SonarQube activo) y eficiencia de recursos para sostenibilidad.

\end{document}
